\documentclass[10pt,twocolumn]{article}
\usepackage[scaled]{helvet}
\renewcommand\familydefault{\sfdefault}
\usepackage[T1]{fontenc}
\usepackage[margin=0.7in]{geometry}
\usepackage{titlesec}
\usepackage{enumitem}
\usepackage{parskip}
\setlength{\columnsep}{0.24in}
\setlist[itemize]{leftmargin=1.2em, topsep=0.2em, itemsep=0.18em}
\titleformat{\section}{\large\bfseries}{\thesection}{1em}{}

\begin{document}
\title{\vspace{-1.6cm}AWS SageMaker Pipelines}
\author{AWS ML Study Notes}
\date{}
\maketitle
\small

\section*{Overview}
SageMaker Pipelines is a managed workflow service for machine learning lifecycle orchestration inside SageMaker. It expresses ML steps as a directed graph with explicit inputs, outputs, parameters, and dependencies.

Unlike generic CI pipelines, it is designed around ML concerns such as processing, training, evaluation, conditional registration, and repeatable data driven execution.

\section*{Core Concepts}
\begin{itemize}
\item Pipelines can include processing steps, training steps, tuning jobs, evaluation logic, conditional branches, cache reuse, and model registration.
\item Pipeline parameters let you vary dataset paths, instance types, thresholds, or model package groups without rewriting the pipeline definition.
\item Step outputs become named references for downstream stages, which improves reproducibility and makes the workflow graph easier to reason about.
\end{itemize}

\section*{How It Works}
\begin{itemize}
\item A pipeline definition is created in code and submitted to SageMaker, where it becomes a reusable managed workflow.
\item When the pipeline runs, preprocessing produces features, training creates a model artifact, evaluation computes quality metrics, and a conditional step decides whether the model should be registered or deployed.
\item Each run leaves an execution record, which helps trace exactly which data, parameters, and images were involved in a given model version.
\end{itemize}

\section*{AI and ML Relevance}
\begin{itemize}
\item This service is useful when retraining or promotion needs to happen repeatedly with discipline, especially across multiple environments or business units.
\item It aligns well with model governance because quality thresholds, approval gates, and lineage are represented as part of the workflow rather than as informal side notes.
\end{itemize}

\section*{Operational Notes}
\begin{itemize}
\item Keep the pipeline definition in version control and treat it like application code. Hidden notebook logic is a common source of unreproducible ML systems.
\item Use caching and incremental design where possible, but ensure cache reuse does not mask a legitimate need to rerun a step after data or code changes.
\item Be explicit about what constitutes a deployable model. A finished training job is not automatically a production ready outcome.
\end{itemize}

\section*{What To Remember}
\begin{itemize}
\item SageMaker Pipelines is about repeatable ML execution, not general source code CI.
\item The core mental model is preprocess, train, evaluate, decide, and promote.
\item A high quality pipeline makes model release criteria visible and enforceable instead of tribal knowledge.
\end{itemize}

\end{document}
